\documentclass[a4paper,11pt]{article}

%--- Pacotes Base --------------------------------------------------------------

\usepackage[portuguese]{babel}
\usepackage[total={6.8in, 10.5in}, top=1cm]{geometry}
\usepackage{ragged2e}

%-------------------------------------------------------------- Pacotes Base ---
%--- Tipografia e Cores --------------------------------------------------------

\usepackage{fontspec}
\setmainfont[
  BoldFont=Cormorant Garamond Bold,
  ItalicFont=Cormorant Garamond Italic,
  BoldItalicFont=Cormorant Garamond Bold Italic
]{Cormorant Garamond}

\usepackage{xcolor}
\definecolor{highlight}{HTML}{283559}

%-------------------------------------------------------- Tipografia e Cores ---
%--- Formatação de Estruturas --------------------------------------------------

\usepackage{enumitem}
\setlist[itemize]{
  noitemsep,
  topsep=2pt,
  parsep=2pt,
  partopsep=0pt,
  leftmargin=15pt
}

\usepackage{titlesec}
\titleformat{\section}{\Large\scshape\raggedright}{}{0em}{}[\titlerule]
\titlespacing*{\section}{0pt}{\baselineskip}{1.5ex}

\pagestyle{empty}

\usepackage[
  colorlinks=true,
  linkcolor=highlight,
  urlcolor=highlight,
  citecolor=highlight
]{hyperref}
\usepackage{bookmark}

%-------------------------------------------------- Formatação de Estruturas ---
%--- Macros Customizadas -------------------------------------------------------

\newcommand{\datedsubsection}[4]{%
  \vspace{0.8ex}
  \noindent\begin{tabular*}{\textwidth}{@{\extracolsep{\fill}} l r@{}}
    \textbf{#3} & \scshape{#2} \\
    \textit{#1} & \\
  \end{tabular*}\par
  \vspace{0.5ex}
  \justifying\noindent
  #4\par
}

% O telefone fica fora do repositório público: o phone.tex, que o git
% ignora, redefine \phone nos builds locais.
\newcommand{\phone}{}
\InputIfFileExists{phone.tex}{}{}

%------------------------------------------------------- Macros Customizadas ---

\begin{document}

%--- Cabeçalho -----------------------------------------------------------------

\begin{center}
  {\Huge Thiago C. Silva} \\
  \vspace*{0.2cm}
  {\scshape Desenvolvedor de Software} \\
  {Curvelo, MG $\cdot$ librefos@newliber.com\phone} \\
  {\href{https://linkedin.com/in/librefos}{linkedin.com/in/librefos} $\cdot$
  \href{https://github.com/librefos}{github.com/librefos}}
\end{center}

%----------------------------------------------------------------- Cabeçalho ---
%--- Resumo Profissional -------------------------------------------------------

\section{Resumo Profissional}
{\setlength{\parindent}{0pt}\setlength{\parskip}{0.5\baselineskip}
Desenvolvedor de software com foco em resolver problemas.  Antes de propor uma
correção, reproduzo o problema e isolo a causa; depois, valido cada etapa com
testes e documento as causas e o porquê de cada decisão, para que ajustes
futuros não criem novos problemas e o conhecimento fique com a equipe, não só
comigo.

Trabalho com sites, sistemas, automações e servidores.  Uso a linha de comando
e sistemas POSIX todos os dias, o que me deixa à vontade para investigar falhas
e comportamentos inesperados direto no sistema.  Minhas automações vão de
scripts em shell e Python a \textit{harnesses} e orquestração de agentes de
LLM.

Sou adepto do software livre e contribuo com o que agrega valor à comunidade.
Sempre em busca de desafios porque quanto maior o desafio, maior a satisfação
e, nesse ciclo, mais eu aprendo.\par}

%------------------------------------------------------- Resumo Profissional ---
%--- Experiência Profissional --------------------------------------------------

\section{Experiência Profissional}

\datedsubsection{outubro de 2026 -- presente}{Coenergy Solar}%
{Desenvolvedor II (Pleno)}{
  \begin{itemize}
    \item Empresa do grupo Cordeiro Energia, fundada em 2025, que oferece
    energia solar por assinatura: o cliente economiza na conta de luz sem
    instalar nada.
    \item Faço parte da equipe que desenvolve e mantém o sistema e a
    infraestrutura da empresa: uma aplicação full stack em React, TypeScript,
    Express e PostgreSQL, rodando em Docker, que atende clientes e pagamentos
    reais.
    \item Faturamento e Pix automáticos, integrados à API do Banco Inter.
    \item Contratos tipados da API, do back-end ao front-end, e testes
    automatizados.
    \item O sistema também roda como aplicativo instalável (PWA).
  \end{itemize}
}

\datedsubsection{julho de 2025 -- presente}{PJ}%
{Desenvolvedor Web Freelancer}{
  \begin{itemize}
    \item \textbf{Flow Digital}, agência de marketing digital (desde maio de
    2026): único desenvolvedor da agência e responsável pelo sucesso da área
    web de todos os clientes.  Desenvolvo, mantenho e atualizo os sites da
    agência e dos clientes, da criação à hospedagem.
    \begin{itemize}
      \item Criei nove sites novos para substituir os antigos (quatro no ar,
      cinco em revisão).  Os quatro no ar têm nota máxima nas avaliações de
      qualidade e velocidade do Google.
      \item Montei a publicação automática dos sites: cada mudança é testada
      antes de ir ao ar.
      \item Automatizei com inteligência artificial tarefas recorrentes da
      agência: relatórios mensais, SEO e calendário de conteúdo.
      \item Criei formulários de contato protegidos contra spam, que enviam
      cada contato direto para o CRM.
      \item Entre os projetos em andamento, um sistema que cria e publica
      posts no Instagram automaticamente: a lógica de base já está pronta, e
      falta polir e integrar à automação principal da agência.
    \end{itemize}
    \item \textbf{Paty Cosméticos}, loja de cosméticos: criei a loja virtual em
    WordPress e WooCommerce, com pagamentos pelo Mercado Pago.
  \end{itemize}
}

\datedsubsection{novembro de 2023 -- julho de 2025}{Flow Digital}%
{Desenvolvedor Web}{
  \begin{itemize}
    \item Cuidei da hospedagem dos sites de cerca de doze clientes:
    servidores, bancos de dados e e-mails; migrações de servidor.
    \item Criei sites e \textit{landing pages}, com foco em SEO.
  \end{itemize}
}

%-------------------------------------------------- Experiência Profissional ---
%--- Projetos Pessoais ---------------------------------------------------------

\section{Projetos Pessoais}
\begin{itemize}
  \item \textbf{Instalação automatizada dos meus computadores}
  (\href{https://github.com/librefos/arch-bootstrap}{arch-bootstrap} e
  \href{https://github.com/librefos/nixos-bootstrap}{nixos-bootstrap},
  GPLv3): do zero, a partir de um repositório, com disco criptografado e
  \textit{Secure Boot}.  Testes em três níveis, inclusive numa máquina virtual
  que faz a instalação completa sem intervenção.
  \item \textbf{Ambiente isolado para agentes de IA}
  (\href{https://github.com/librefos/corral}{corral}, GPLv3): contêiner
  descartável e sem privilégios para agentes de programação (Claude Code,
  Codex, opencode e agy), com sistema de arquivos somente leitura,
  \textit{backup} criptografado e um \textit{harness} próprio que cria,
  atualiza e destrói o ambiente.
  \item \textbf{Central de CRM e atendimento}
  (\href{https://github.com/librefos/tcrm}{tcrm}, GPLv3): CRM, chat de
  atendimento, WhatsApp e automações (n8n) rodando juntos em Docker, atrás de
  um proxy reverso e com um banco de dados por aplicação.  O CRM fica aberto a
  agentes de IA via MCP.
  \item \textbf{Servidores próprios}: um servidor Git na minha rede local e um
  servidor virtual, acessado por VPN, com e-mail próprio no meu domínio (SPF,
  DKIM e DMARC) e IRC sempre conectado.
  \item \textbf{Currículo vivo}
  (\href{https://github.com/librefos/resume}{resume}, GPLv3): este currículo
  é um projeto em \LaTeX{}, atualizado conforme a minha área e a minha
  experiência evoluem, e versionado no Git como todo software que faço.  A
  cada nova versão, uma integração contínua própria gera os PDFs em português
  e em inglês num contêiner e os publica para download.
\end{itemize}

%--------------------------------------------------------- Projetos Pessoais ---
%--- Contribuições em Código Aberto --------------------------------------------

\section{Contribuições em Código Aberto}
\begin{itemize}
  \item \href{https://patchwork.sourceware.org/project/binutils/list/?submitter=37677}%
  {GNU Binutils}, ferramentas do projeto GNU usadas para compilar o Linux:
  duas correções aceitas (2026).
  \item \href{https://github.com/lichess-org/lila/pull/21832}{Lichess},
  servidor de xadrez on-line livre e gratuito: uma correção aceita (2026).
\end{itemize}

%-------------------------------------------- Contribuições em Código Aberto ---
%--- Habilidades ---------------------------------------------------------------

\section{Habilidades}
\begin{itemize}
  \item \textbf{Programação:} Shell, Python, JavaScript, TypeScript, PHP, Java,
  C e C++.
  \item \textbf{Web:} React, Express, Astro, PWA, SEO, GEO e integração com
  APIs e MCP.
  \item \textbf{Testes e CI/CD:} Vitest, Playwright, Lighthouse CI, Git e
  GitHub Actions.
  \item \textbf{Infraestrutura:} Linux, VPS, hospedagem compartilhada, Apache,
  Nginx, Docker, Podman e Nix.
  \item \textbf{Bancos de dados:} MySQL e PostgreSQL.
  \item \textbf{Documentação:} \LaTeX{}, Texinfo e Typst.
  \item \textbf{IA:} agentes LLM (Claude Code, Codex, opencode, agy, pi) e
  orquestração de agentes.
  \item \textbf{Idiomas:} português nativo e inglês intermediário.
\end{itemize}

%--------------------------------------------------------------- Habilidades ---
%--- Formação e Certificações --------------------------------------------------

\section{Formação Acadêmica \& Certificações}
\textbf{Bacharelado em Engenharia de Software} $\cdot$ UniCesumar
\hfill 2023 -- 2027 \\
\textbf{Técnico em Desenvolvimento de Sistemas} $\cdot$ Senac Minas
\hfill 2023 -- 2024 \\
\textbf{GitHub Foundations} $\cdot$ GitHub \hfill 2024 \\
\textbf{Cybersecurity Essentials} $\cdot$ Linux Foundation \hfill 2024

%-------------------------------------------------- Formação e Certificações ---

\end{document}
